\documentclass[10pt,a4paper]{amsart}
\usepackage[margin=19mm]{geometry}
\usepackage{amsmath,amssymb,mathtools}
\usepackage[scr=rsfs,cal=euler]{mathalfa}
\usepackage{xfrac}
\usepackage[unicode,hidelinks,bookmarks=false]{hyperref}
\newcommand{\ie}{i.e.\ }
\newcommand{\cf}{cf.\ }
\newcommand{\resp}{respectively\ }
\newcommand{\Subsection}{\subsection}
\providecommand{\nobreakdash}{\nobreak\hbox{-}}
\setlength{\emergencystretch}{2em}
\multlinegap0pt
\usepackage{bm}
\usepackage{bbm}
\usepackage{dsfont}
\usepackage{xspace}
\usepackage{cancel}
\usepackage{mathtools}
\usepackage{amsthm}
\makeatletter
\@ifundefined{theorem}{\newtheorem{theorem}{Theorem}}{}
\@ifundefined{proposition}{\newtheorem{proposition}[theorem]{Proposition}}{}
\makeatother
\title[Memory-Type Null Controllability of Heat Equations with Dela]{Memory-Type Null Controllability of Heat Equations with Delay Effects}
\author{Dev Prakash Jha, Raju K. George}
\date{Source: arXiv:2506.22221v1 (CC BY 4.0)}
\begin{document}
\maketitle

\noindent\emph{Source and licence.} Memory-Type Null Controllability of Heat Equations with Delay Effects. Version arXiv:2506.22221v1 (CC BY 4.0), distributed under the Creative Commons Attribution 4.0 International licence, as recorded in the arXiv licence metadata for this exact version and shown on the abstract page https://arxiv.org/abs/2506.22221v1. Full rights notice: licenses/arxiv-2506.22221-CC-BY-4.0.txt.\par\smallskip

\noindent\textit{Editorial scope.} Counted unit: the Theorem whose source label is \texttt{th:Finite\_the\_1} in the article (source0.tex lines 539--552, source label \texttt{th:Finite\_the\_1}), delivered together with the article's own complete original proof of it (source0.tex lines 555--688, one \texttt{proof} environment of 134 lines). No mathematics is written, repaired or re-derived: statement and proof are the article's own text line for line. Required context retained uncounted: eq:Intro\_eq\_7 (lines 210--218); eq:dual\_eq\_1 (lines 337--346); prop:2.1 (lines 354--362); eq:Finite\_eq\_1 (lines 511--518); eq:Finite\_eq\_2 (lines 524--533). Every definition, display and label of those lines is retained unchanged. Omitted from this package and not redistributed: 1--209, 219--336, 347--353, 363--510, 519--523, 534--538, 553--554, 689--1814. Nothing inside a retained excerpt is omitted or altered: every retained body is delivered exactly as the article's lines are written, so no omission receipt of any kind accompanies this package. Citations: the retained excerpts carry no citation command at all, so no bibliography entry of the article is transcribed and no reference is redistributed. Reference adaptation: none was needed and none was made; every reference inside the delivered unit resolves inside it, so no unresolved reference or citation is produced. Printed slips kept literal: no printed word, symbol, hypothesis or number of the counted statement or of its proof was altered. Layout and class: the article's own class is \texttt{article} with packages including lipsum, hyperref, authblk, geometry, graphicx, multirow, amsmath, amssymb, amsfonts, amsthm, mathrsfs, appendix, xcolor, textcomp; the wrapper is the lane's pinned \texttt{amsart} editorial wrapper at 10pt on A4, which restates the theorem-like environments the retained text needs and adds the article's own macro definitions that the retained text needs (none), with extra wrapper packages (bm, bbm, dsfont, xspace, cancel, mathtools). The delivered numbering is the unit's own. Assets: no retained span contains a figure environment, a graphics-inclusion command, a PDF-page inclusion, a table or a TikZ picture; no image file is copied into this package, the package declares no asset dependency and no image file is redistributed with it. Licence: version arXiv:2506.22221v1 of this article is distributed under the Creative Commons Attribution 4.0 International licence, as recorded in the arXiv licence metadata for this exact version and shown on the abstract page https://arxiv.org/abs/2506.22221v1. A full rights notice is delivered at licenses/arxiv-2506.22221-CC-BY-4.0.txt. Characters: every retained excerpt is pure ASCII, so the delivered engine, XeLaTeX with its default Latin Modern text font, reports no missing character.
\begin{equation}
\label{eq:Intro_eq_7}
\left\{
\begin{aligned}
    y_t(t) &= Ay(t) + A_1 y(t-h) + \int_0^t M(t - s)y(s)\,ds + B(t)u(t), \quad t \in [0, T], \\
    y_0(\theta) &= \varphi(\theta),\quad \theta \in [-h, 0].
\end{aligned}
\right.
\end{equation}

\begin{equation}
\label{eq:dual_eq_1}
\left\{
\begin{aligned}
    -w_t(t) &= A^* w(t) + A_1^* w(t + h) + \int_t^T M(s - t)^* w(s)\,ds - \widetilde{M}(T - t)^* z_T, \quad t \in [0, T], \\
    w(T) &= w_T, \\
    w(t) &= 0, \quad t \in (T, T + h].
\end{aligned}
\right.
\end{equation}

\begin{proposition}
\label{prop:2.1}
The system described by \eqref{eq:Intro_eq_7} is Delay and memory-type null controllable (with the kernel $\widetilde{M}(\cdot)$) if and only if every solution to \eqref{eq:dual_eq_1} satisfies the observability inequality:
\begin{equation}
\label{eq:enhanced_obs}
\|w(\theta)\|_Y^2 + \int_{ - h}^{\theta} \|w(t+h)\|_Y^2 dt+ \int_{T_1 - h}^{T-h} \|w(t+h)\|_Y^2 dt \leq K \int_0^T \|B(t)^* w(t)\|_U^2 dt,
\end{equation}
for all $w_T, z_T \in Y$ and $T_1 \in [T - h, T]$
\end{proposition}

\begin{equation}\label{eq:Finite_eq_1}
\left\{
\begin{aligned}
y_t &= Ay(t)+A_1 y(t-h) + \int_0^t M(t-s)y(s)\,ds + Bu(t), \quad t \in (0,T], \\
y(\theta) &= \varphi (\theta), \quad \theta \in [-h, 0].
\end{aligned}
\right.
\end{equation}

\begin{equation}
\label{eq:Finite_eq_2}
\left\{
\begin{aligned}
    w_t(t) &= -A^* w(t) - A_1^* w(t + h) - \int_t^T M(s - t)^* w(s)\,ds +\widetilde{M}(T - t)^* z_T, \quad t \in [0, T], \\
    w(T) &= w_T, \\
    w(t) &= 0, \quad t \in (T, T + h].
\end{aligned}
\right.
\end{equation}

\begin{theorem}\label{th:Finite_the_1}
    \begin{itemize}
    \item[(i)] Suppose that \( M(\cdot), \widetilde{M}(\cdot) \in L^1(0,T;\mathbb{R}^{n \times n}) \). Then for any solution \( w \) of \eqref{eq:Finite_eq_2}, the condition
    \begin{equation}\label{eq:Finite_eq_3}
         B^T w = 0 \text{ in } [0,T] \quad \Rightarrow \quad x_T = \widetilde{M}(t)^T x_T = 0 \quad \text{for almost every } t \in [0,T]
    \end{equation}
    ensures that the system \eqref{eq:Finite_eq_1} is Delay and memory-type null controllable.

    \item[(ii)] Assume \( M(\cdot) = G M'(\cdot) \) and \( \widetilde{M}(\cdot) = \widetilde{G} M'(\cdot) \) for some constant matrices \( G, \widetilde{G} \in \mathbb{R}^{n \times n} \). If the system \eqref{eq:Finite_eq_1} is Delay and memory-type null controllable, then for every solution of \eqref{eq:Finite_eq_2}, it holds that
     \begin{equation}\label{eq:Finite_eq_4}
         B^T w = 0 \text{ in } [0,T] \quad \Rightarrow \quad \widetilde{M}(t)^T x_T = 0 \quad \text{for almost every } t \in [0,T].
    \end{equation}
    \end{itemize}
\end{theorem}

\begin{proof}
To prove (i), by  using (\ref{eq:Finite_eq_3}), we devide in multipal steps: \\
We aim to prove the inequality:
\begin{equation}\label{eq:obs}
|w_T|^2 + \left( \int_0^T |\widetilde{M}(t)^\top z_T|\, dt \right)^2 \leq C \int_0^T |B^\top w(t)|^2\, dt,
\end{equation}
under the assumption that for any solution $w(t)$ to the adjoint system,
\[
B^\top w(t) \equiv 0 \text{ on } [0, T] \Rightarrow w_T = 0 \text{ and } \widetilde{M}(t)^\top z_T = 0 \text{ a.e. } t \in [0, T].
\]

\textit{Step 1: Contradiction Hypothesis}

Assume, for contradiction, that inequality \eqref{eq:obs} fails. Then for every integer $k \in \mathbb{N}$, there exist terminal data $(w_T^k, z_T^k) \in \mathbb{R}^n \times \mathbb{R}^n$ and corresponding solution $w^k(t)$ of the adjoint system such that:
\[
\int_0^T |B^\top w^k(t)|^2 dt < \frac{1}{k^2} \left( |w_T^k|^2 + \left( \int_0^T |\widetilde{M}(t)^\top z_T^k| dt \right)^2 \right).
\]

\textit{Step 2: Normalize and Pass to the Limit}

Define normalization constants:
\[
N_k := \sqrt{ |w_T^k|^2 + \left( \int_0^T |\widetilde{M}(t)^\top z_T^k| dt \right)^2 } > 0,
\]
and set:
\[
\hat{w}_T^k = \frac{w_T^k}{N_k}, \quad \hat{z}_T^k = \frac{z_T^k}{N_k}.
\]

Let $\hat{w}^k(t)$ denote the solution to the adjoint system with terminal condition $(\hat{w}_T^k, \hat{z}_T^k)$. Then:
\[
\int_0^T |B^\top \hat{w}^k(t)|^2 dt < \frac{1}{k^2}, \quad \text{and} \quad |\hat{w}_T^k|^2 + \left( \int_0^T |\widetilde{M}(t)^\top \hat{z}_T^k| dt \right)^2 = 1.
\]

By compactness in finite dimensions, we may extract a weakly converging subsequence:
\[
\hat{w}_T^k \rightharpoonup w_T^*, \quad \hat{z}_T^k \rightharpoonup z_T^*.
\]

Moreover, the corresponding solutions satisfy:
\[
B^\top \hat{w}^k(t) \to 0 \text{ in } L^2(0,T) \Rightarrow B^\top w^*(t) = 0.
\]

\textit{Step 3: Use the Hypothesis}

By the assumption, $B^\top w^*(t) \equiv 0$ implies:
\[
w_T^* = 0, \quad \widetilde{M}(t)^\top z_T^* = 0 \text{ for all } t.
\]

Thus:
\[
|w_T^*|^2 + \left( \int_0^T |\widetilde{M}(t)^\top z_T^*| dt \right)^2 = 0.
\]

\textit{Step 4: Contradiction}

But by weak lower semi-continuity of the norm:
\[
\liminf_{k \to \infty} \left( |\hat{w}_T^k|^2 + \left( \int_0^T |\widetilde{M}(t)^\top \hat{z}_T^k| dt \right)^2 \right) \geq 1,
\]
contradicting the convergence to zero of the limit. Therefore, the inequality \eqref{eq:obs} must hold.\\


\textit{Step 5: Energy estimation }

We aim to prove the estimate:
\begin{equation}
\label{eq:energy_estimate}
\|w(\theta)\|_Y^2 + \int_{ - h}^{\theta} \|w(t+h)\|_Y^2 dt+ \int_{T_1 - h}^{T-h} \|w(t+h)\|_Y^2 dt \leq C \left[ |w_T|^2_{\mathbb{R}^n} + \left( \int_0^T \widetilde{M}(t)^\top z_T \, dt \right)^2 \right] \quad \forall w_T, z_T \in \mathbb{R}^n.
\end{equation}

We do not attempt a direct energy computation due to the presence of a future-integral memory term in the adjoint system:
\begin{equation}
\label{eq:adjoint}
\dot{w}(t) = -A^\top w(t) - \int_t^T M(s - t)^\top w(s) \, ds + \widetilde{M}(T - t)^\top z_T, \quad w(T) = w_T.
\end{equation}

Instead, we use the well-posedness and stability theory for linear integro-differential equations. From standard results on such systems (see, e.g., Dafermos or Lions), there exists a constant $C > 0$ such that the solution $w(\cdot)$ satisfies the stability estimate:
\[
\|w\|_{C([0, T]; \mathbb{R}^n)} \leq C \left( |w_T|_{\mathbb{R}^n} + \left\| \widetilde{M}(\cdot)^\top z_T \right\|_{L^1(0, T; \mathbb{R}^n)} \right).
\]

In particular, evaluating at $t = 0$ gives:
\[
|w(0)|_{\mathbb{R}^n} \leq C \left( |w_T|_{\mathbb{R}^n} + \int_0^T |\widetilde{M}(t)^\top z_T| \, dt \right).
\]

Applying the Cauchy--Schwarz inequality, we estimate:
\[
\left( \int_0^T |\widetilde{M}(t)^\top z_T| \, dt \right)^2 \leq T \int_0^T |\widetilde{M}(t)^\top z_T|^2 \, dt \leq T \| \widetilde{M} \|_{L^2(0, T)}^2 \cdot |z_T|^2.
\]

Combining the estimates and absorbing constants into $C$, we obtain:
\[
\|w(\theta)\|_Y^2 + \int_{ - h}^{\theta} \|w(t+h)\|_Y^2 dt+ \int_{T_1 - h}^{T-h} \|w(t+h)\|_Y^2 dt \leq C \left( |w_T|^2_{\mathbb{R}^n} + \left( \int_0^T \widetilde{M}(t)^\top z_T \, dt \right)^2 \right),
\]
which completes the proof of inequality \eqref{eq:energy_estimate}.
Combining (\ref{eq:obs}) and (\ref{eq:energy_estimate}), we show that
\begin{equation}\label{eq:Finite_eq_6}
    \|w(\theta)\|_Y^2 + \int_{ - h}^{\theta} \|w(t+h)\|_Y^2 dt+ \int_{T_1 - h}^{T-h} \|w(t+h)\|_Y^2 dt \leq C\int_0^T |B^{T}w(t)|_{\mathbb{R}^m}^2dt \quad \forall w_T,z_T \in \mathbb{R}^n
\end{equation}


Next, we address (ii). From Proposition \ref{prop:2.1} and the null controllability of (\ref{eq:Finite_eq_1}), we conclude that any solution of (\ref{eq:Finite_eq_2}) satisfies the energy bound in (\ref{eq:Finite_eq_6}). Thus, \( w(\theta) = 0 \). Since \( B^T w \equiv 0 \) on \( [0,T] \), define \( \varphi = w_t \). Referring to the first equation in (\ref{eq:Finite_eq_2}) and using \( M(\cdot) = G M'(\cdot) \), \( \widetilde{M}(\cdot) = \widetilde{G} M'(\cdot) \), we obtain:
\begin{align*}
\varphi_t &= -A^T \varphi(t)-A_1^T \varphi(t+h) + M(0)^T w + \int_t^T M(s - t)^T w(s)ds - \widetilde{M}(T - t)^T x_T \\
&= -A^T \varphi (t)-A_1^T \varphi(t+h) + M(0)^T w + \int_t^T M(s - t)^T w(s)ds - \widetilde{M}(T - t)^T x_T \\
&= -A^T \varphi(t)-A_1^T \varphi(t+h) - \int_t^T M(s - t)^T \varphi(s)ds + \widetilde{M}(T - t)^T (G^T w_T - \widetilde{G}^T x_T).
\end{align*}

Thus, \( \varphi \) satisfies:
\begin{equation}
\label{eq:Finite_eq_8}
\left\{
\begin{aligned}
\varphi_t &= -A^T \varphi(t)-A_1^T \varphi(t-h) - \int_t^T M(s - t)^T \varphi(s)ds + \widetilde{M}(T - t)^T (G^T w_T - \widetilde{G}^T x_T), \quad t \in [0, T), \\
\varphi(T) &= -A^T w_T-A^T w(T-h) + \widetilde{M}(0)^T z_T.
\end{aligned}
\right.
\end{equation}

Recognizing that (\ref{eq:Finite_eq_8}) shares the same structure as (\ref{eq:Finite_eq_2}), the energy estimate (\ref{eq:Finite_eq_6}) implies
\begin{equation}
\label{eq:Finite_eq_9}
 \|\varphi(\theta)\|_Y^2 + \int_{ - h}^{\theta} \|\varphi(t+h)\|_Y^2 dt+ \int_{T_1 - h}^{T-h} \|\varphi(t+h)\|_Y^2 dt \leq C \int_0^T |B^T \varphi(s)|^2 ds = C \int_0^T |B^T w_t(s)|^2 ds = 0.
\end{equation}

Consequently, \( w_t(\theta) = 0 \). Iterating this argument shows that all time derivatives vanish at \( t = \theta \), i.e., \( \frac{d^k w(t)}{dt^k}\big|_{t=\theta} = 0 \) for all \( k \in \mathbb{N}_0 \). Since \( w(t) \) is analytic in time, this leads to \( w(t) \equiv 0 \) on \( [0,T] \). Hence,
\[
w_T = \widetilde{M}(t)^T z_T = 0 \quad \text{for all } t \in [0, T].
\]
\end{proof}

\end{document}
